\begin{abstract}
Retrieval-augmented speech translation improves rare terminology, yet systems
such as Retrieval-Augmented Simultaneous Speech Translation, or RASST, expect
the user to supply a glossary before every session. Merging every domain into
one universal glossary removes this step, but as the catalog expands,
irrelevant and conflicting entries crowd out useful terms inside the fixed
retrieval and prompt budget. We present AutoTerm-SST, an interactive streaming
speech translation system that operates without a glossary upload or domain
selection. AutoTerm-SST partitions an offline terminology catalog into topic
slices, matches recent partial translations against a short bilingual
description of each slice, and restricts term retrieval to the best-matching
slices under an unchanged budget. On a 4.7-hour English-to-Chinese stream
alternating five ACL talks with five oncology lectures, AutoTerm-SST reaches
88.3\% terminology accuracy, showing no significant difference from a
known-domain reference at 88.1\%, and a second independent run reproduces this
parity. Against a one-million-entry merged glossary, AutoTerm-SST gains 3.9
points in terminology accuracy and 2.30 BLEU, both significant in a talk-level
permutation test. The full pipeline runs 2.46 times faster than real time, and
the live interface displays active slices and retrieved term pairs alongside
every partial translation.
\end{abstract}
